\chapter{EPS – Sistema de Energía}

\section{Visión general del Subsistema}
El subsistema de Energía Eléctrica (EPS, \textit{Electrical Power System}) es responsable primario de tres funciones en órbita LEO: (i) recolección energética por paneles solares, (ii) almacenamiento mediante celdas de batería y (iii) acondicionamiento y distribución segura a todo el bus de carga. Adicionalmente, el EPS orquesta indirectamente la vida del satélite al controlar los interruptores principales (Kill-Switches).

\section{Generación de Energía (Paneles Solares)}
La fuente de energía para INTISAT está dotada por arreglos adosados de celdas solares (típicamente celdas fotovoltaicas de triple unión de alta eficiencia y certificación espacial, como Emcore o AzurSpace) dispuestas a lo largo y ancho de caras externas del prisma del diseño CubeSat.
\begin{itemize}
    \item \textbf{Gestión de Generación:} La tarjeta del procesador del EPS procesa la entrada de energía a través de algoritmos lógicos de Seguimiento de Punto de Máxima Potencia (MPPT - Maximum Power Point Tracking), garantizando maximizar el flujo en condiciones muy variables de irradiación térmica y lumínica y optimizando la eficiencia de los paneles.
\end{itemize}

\section{Almacenamiento y Capacidad (Battery Module)}
Todo exceso de energía cosechada respecto a los consumos inmediatos en la fase soleada (In-Sunlight) se direcciona de manera prioritaria al \textit{Battery Module}. 
\begin{itemize}
    \item \textbf{Celdas:} compuesto por celdas robustas de iones de litio integradas en serie/paralelo (formato 18650), con calentadores térmicos resistivos para sobrepasar los eclipses orbitales LEO sin congelar los electrolitos.
    \item \textbf{Battery Management System (BMS):} Mide constantemente la profundidad de descarga (Depth of Discharge - DoD), voltajes diferenciales de las celdas, amperajes y termistores, realimentando al módulo EPS con parámetros de advertencia para evitar sub-voltajes destructivos.
\end{itemize}

\section{Acondicionamiento, Distribución y Seguridad}
La potencia es segmentada de la batería para transformarla mediante sistemas estabilizadores (DC-DC) a líneas reguladas estándares (ej. 3.3V y 5V) y la línea nativa no regulada ($V_{batt}$).
\begin{itemize}
    \item \textbf{Protección Analógica e Interfaz de Bus:} Provee una protección inmediata de hardware mediante mitigación de Latch-Ups Locales (LCL) que detectan cortocircuitos temporales inducidos por radiación o consumo desbordado de las cargas útiles cortando la energía de puertos individuales.
    \item \textbf{Integración (Interface Data Sheet):} El EPS se comunica de manera bidireccional mediante el I2C/CAN cruzado del PC-104 reportando telemetría de sanidad a la OBC e indicando el porcentaje de la batería constante, accionador central de las rutinas contingentes.
\end{itemize}
